\documentclass[11pt,letterpaper]{article}

% --- Paquetes de Idioma y Tipografía ---
\usepackage[utf8]{inputenc}
\usepackage[spanish]{babel}
\usepackage[T1]{fontenc}
\usepackage{lmodern}
\usepackage{microtype}

% --- Configuración de Página y Márgenes ---
\usepackage[top=2.0cm, bottom=2.0cm, left=2.2cm, right=2.2cm]{geometry}
\usepackage{setspace}
\setstretch{1.12}

% --- Paquetes Matemáticos y Cajas ---
\usepackage{amsmath, amssymb}
\usepackage[most]{tcolorbox}

% --- Enlaces ---
\usepackage{hyperref}
\hypersetup{colorlinks=true, linkcolor=blue!70!black, urlcolor=blue!70!black}

% --- Encabezados ---
\usepackage{fancyhdr}
\pagestyle{fancy}
\fancyhf{}
\fancyhead[L]{\small Computación Científica --- Primera Entrega}
\fancyhead[R]{\small Guía Rápida de Sustentación Individual (12 min)}
\fancyfoot[C]{\thepage}
\renewcommand{\headrulewidth}{0.4pt}

\newtcolorbox{pregunta}[2][]{
    colback=blue!4!white,
    colframe=blue!70!black,
    fonttitle=\bfseries,
    title=#2,
    #1
}

\begin{document}

\begin{center}
    {\Large \textbf{Guía Maestra de Preparación para la Sustentación Individual}}\\[0.15cm]
    {\large \textbf{Primera Entrega de Proyecto: Exponente de Lyapunov en Mapas 1D}}\\[0.1cm]
    \small \textbf{Modalidad:} 12 minutos por grupo. Preguntas individuales directas del profesor Juan Fernando Paz. \\
    \textbf{Regla de oro:} Responder de forma rápida, simple, intuitiva y conceptualmente sólida.
\end{center}
\vspace{0.2cm}

\begin{pregunta}{1. En palabras muy sencillas, ¿qué es el caos y qué mide el Exponente de Lyapunov?}
\textbf{Respuesta intuitiva (Modelo mental para el póster):}
\begin{itemize}
    \item \textbf{Caos:} No es azar ni ruido; es una fórmula matemática determinista exacta cuya trayectoria se vuelve impredecible porque cualquier error o diferencia inicial microscópica crece exponencialmente.
    \item \textbf{Exponente de Lyapunov ($\lambda$):} Es el «termómetro cuantitativo» del caos. Si imaginamos dos puntos casi idénticos separados por $\delta_0$, en cada iteración la distancia se multiplica por un factor de estiramiento (la derivada local $|f'(x)|$). $\lambda$ mide la tasa promedio de ese estiramiento:
    \begin{itemize}
        \item Si $\lambda < 0$: Las trayectorias se atraen y se comprimen hacia un punto fijo (predecible).
        \item Si $\lambda = 0$: Neutro, separación lineal (borde de transición).
        \item Si $\lambda > 0$: Caos puro. Las trayectorias se separan exponencialmente ($e^{\lambda n}$). Entre mayor sea $\lambda$, más rápido se pierde la capacidad de predecir el futuro.
    \end{itemize}
\end{itemize}
\end{pregunta}

\begin{pregunta}{2. ¿Cómo se conecta este proyecto directamente con lo visto en las clases del curso?}
\textbf{Respuesta clave (Puntos del programa de clase):}
\begin{enumerate}
    \item \textbf{Método del Punto Fijo (\textit{Clase: Punto Fijo}):} En clase vimos la iteración $x_{n+1} = g(x_n)$ y demostramos con el Teorema del Valor Medio que el error evoluciona como $e_{n+1} \approx g'(\xi_n) e_n$. Si $|g'(x^*)| < 1$, el método converge. En un mapa caótico, se tiene iteración sucesiva $x_{n+1} = f(x_n)$ pero con un promedio $|f'(x)| > 1$. ¡El exponente de Lyapunov $\lambda = \frac{1}{N} \sum \ln |f'(x_i)|$ es literalmente el promedio del factor de amplificación de error que analizamos en la iteración de punto fijo!
    \item \textbf{Aritmética IEEE 754 y Épsilon de Máquina (\textit{Clase 2}):} En clase estudiamos $\varepsilon_{\text{mach}} \approx 2.22 \times 10^{-16}$. En un sistema caótico ($\lambda > 0$), este error de redondeo inicial se amplifica por $e^{\lambda n}$, haciendo que dos simulaciones en computadoras con ligeras diferencias diverjan totalmente tras unas 50 iteraciones (degradación digital).
    \item \textbf{Valores Propios y Singulares (\textit{Clases: Autovalores y SVD}):} Los valores singulares miden cuánto estira o contrae una transformación matricial el espacio. En mapas 1D, la derivada $|f'(x)|$ es el factor de estiramiento instantáneo y $\lambda$ es su tasa asintótica acumulada.
\end{enumerate}
\end{pregunta}

\begin{pregunta}{3. ¿Por qué el Exponente de Lyapunov es el único método para verificar qué tan caótico es un sistema?}
\textbf{Respuesta clave:}  
Porque la simple inspección visual de una gráfica o de una serie de números es engañosa: un sistema puede verse desordenado pero ser simplemente un ciclo periódico muy largo. El exponente de Lyapunov es la única métrica matemática rigurosa que calcula la tasa exponencial de divergencia entre órbitas vecinas. Si $\lambda > 0$, matemáticamente existe hipersensibilidad a las condiciones iniciales, lo que certifica la presencia indiscutible de caos determinista.
\end{pregunta}

\begin{pregunta}{4. ¿Cómo se calcularon las soluciones analíticas exactas para validar el código?}
\textbf{Respuesta clave:}
\begin{itemize}
    \item \textbf{Mapa de Tienda:} La pendiente es $|f'(x)| = p$ casi en todas partes ($x \neq 0.5$). Al aplicar la fórmula $\lambda = \frac{1}{N} \sum \ln |f'(x_i)| = \frac{1}{N} \sum \ln(p) = \ln(p)$ de forma exacta.
    \item \textbf{Mapa Logístico en Punto Fijo ($1 < r \le 3$):} El atractor es $x^* = 1 - 1/r$. Al evaluar la derivada $f'(x^*) = r(1 - 2x^*) = 2 - r$, obtenemos $\lambda = \ln|2 - r|$.
    \item \textbf{Mapa Logístico en Caos Total ($r=4$):} Integrando contra la densidad de probabilidad $\rho(x) = \frac{1}{\pi\sqrt{x(1-x)}}$, el valor exacto es $\ln(2) \approx 0.693147$.
\end{itemize}
\end{pregunta}

\begin{pregunta}{5. ¿Qué diferencias principales observaron al graficar los espectros de ambos mapas?}
\textbf{Respuesta clave:}  
\begin{itemize}
    \item El \textbf{Mapa Logístico} tiene un espectro muy irregular: pasa de punto fijo ($\lambda < 0$) a duplicaciones de período y finalmente a caos, pero en medio del caos sufre caídas abruptas a $\lambda < 0$ («ventanas periódicas», como la ventana de período 3 de Li-Yorke en $r \approx 3.83$).
    \item El \textbf{Mapa de Tienda} tiene un espectro suave, continuo y monótonamente creciente $\lambda(p) = \ln(p)$. No tiene ninguna ventana periódica. Para cualquier $p > 1$, el caos está 100\% garantizado.
\end{itemize}
\end{pregunta}

\begin{pregunta}{6. ¿Cómo implementaron el algoritmo en Python y cómo evitaron fallos numéricos?}
\textbf{Respuesta clave:}  
En el módulo \texttt{src/lyapunov.py}:
\begin{enumerate}
    \item \textbf{Descarte de Transitorio:} Descartamos las primeras 500 iteraciones sin sumar derivadas, para que el número inicial se «olvide» y la trayectoria caiga en el atractor real.
    \item \textbf{Salvaguarda contra $\ln(0)$:} Si la derivada pasa cerca de cero ($|f'(x)| < 10^{-14}$), usamos $\max(|f'(x)|, 10^{-14})$ para evitar indeterminaciones por logaritmo de cero.
    \item \textbf{Suma Logarítmica:} Sumamos logaritmos en vez de multiplicar derivadas para evitar que el número desborde (\textit{overflow}) o colapse a cero (\textit{underflow}).
\end{enumerate}
\end{pregunta}

\begin{pregunta}{7. ¿Qué es la «degradación digital» en el mapa de tienda con $p=2$?}
\textbf{Respuesta clave:}  
En computadoras con estándar IEEE 754 de 64 bits, multiplicar por $p=2$ desplaza los bits de la mantisa binaria hacia la izquierda. Si el valor $x$ es una fracción binaria exacta, tras unas 50 iteraciones la mantisa se vacía de unos y se convierte exactamente en $0.0$, deteniendo el caos. En nuestro algoritmo lo resolvimos añadiendo una micro-perturbación de orden $10^{-14}$ si el número cae exactamente en $0.0$ o $0.5$.
\end{pregunta}

\begin{pregunta}{8. ¿Cuáles son los siguientes pasos hacia el póster final (Segunda Entrega)?}
\textbf{Respuesta clave:}  
\begin{enumerate}
    \item Diseñar visualmente el póster sintetizando la analogía del estiramiento, los espectros gráficos de Lyapunov y el glosario para el público.
    \item Usar el mapa de tienda como generador de números pseudoaleatorios para cifrar una imagen mediante permutación y difusión (arquitectura de Fridrich).
    \item Validar el cifrado con las pruebas de seguridad estadística NPCR ($>99.6\%$) y UACI ($\approx 33.4\%$) de Zhang \& Liu (2023).
\end{enumerate}
\end{pregunta}

\end{document}
